\section{Conclusion and Future Work}\label{sec:conclusion}

A hybrid SLC/MLC layout is necessary for SVD-factorised transformers on ReRAM, and which components are
protected matters as much as how many. Ranking by singular value is close to the best placement under a
cell-level error model and needs no data; a small SLC share then suffices at moderate noise, and the hybrid
stays robust to read noise, level-dependent noise, a 6-bit ADC and moderate IR drop. Stuck cells are the
exception. Mapping components to slots around a tested chip's faults recovers most of the loss at no
hardware cost, and tier-aware training closes the remaining gap to an all-SLC chip up to 1\% scattered stuck
cells, with a third fewer cells.

Future work: energy, area and latency measured by emulating the tiers and their periphery on an FPGA board
(unavailable to us during this project) or with NeuroSim~\cite{peng2019neurosim}; all error sources combined
on one chip; fine-tuning and training repeated over several seeds, which our free-tier GPU budget did not
allow; finer activation scales; robustness to fault-injection attacks such as the bit-flip
attack~\cite{rakin2019bfa}; and further tasks and models.
